%% Appendix: Artifact Description / Artifact Evaluation
%% Include from main.tex via %% Appendix: Artifact Description / Artifact Evaluation
%% Include from main.tex via %% Appendix: Artifact Description / Artifact Evaluation
%% Include from main.tex via %% Appendix: Artifact Description / Artifact Evaluation
%% Include from main.tex via \input{CloverIS/paper/appendix.tex}

\section*{Appendix: Artifact Description/Artifact Evaluation}

\subsection*{Artifact Description (AD)}

\subsubsection*{I. Overview of Contributions and Artifacts}

\paragraph{A. Paper's Main Contributions}
The paper introduces \textsf{Clover}, a parallel $k$-clique counting
algorithm based on a Cover-oriented pivoting search, and its extension
\textsf{Clover+IS}, which counts independent sets on the complement graph
when a residual subproblem becomes dense. In the context of the artifact
description, the relevant contributions are:

\begin{description}
    \item[$C_1$] \textsf{Clover} and \textsf{Clover+IS} outperform the
    state-of-the-art \textsf{PivotScale} baseline by one to four orders of
    magnitude for large cliques on large real-world graphs.
    \item[$C_2$] \textsf{Clover+IS} is the first system to count
    $k$-cliques of arbitrary $k$ on \textsf{com-LiveJournal}, a problem
    that was previously intractable.
\end{description}

\paragraph{B. Computational Artifacts}
\begin{center}
\begin{tabular}{lll}
\toprule
Artifact ID & Contributions & Related Paper Elements \\
\midrule
$A_1$ & $C_1$, $C_2$ & Figure~5, Table~II \\
\bottomrule
\end{tabular}
\end{center}

$A_1$: \url{https://github.com/<TBD>/CloverIS} (final DOI assigned via
Zenodo by the Artifact Freeze deadline).

\subsubsection*{II. Artifact Identification}

\paragraph{A. Computational Artifact $A_1$}

\noindent\emph{Relation to contributions.}
$A_1$ contains the full source of \textsf{Clover}, \textsf{Clover+IS}, and
the \textsf{PivotScale} baseline, together with the download/convert
pipeline for the five input graphs and the plotting driver. Running the
single entry-point script produces the per-$(graph, k)$ counting times
that substantiate $C_1$ and $C_2$.

\noindent\emph{Expected results.}
The faster of \textsf{Clover} and \textsf{Clover+IS} outperforms
\textsf{PivotScale} on every graph at every~$k$ tested.
\textsf{Clover+IS} is fastest for $k \geq 5$; at smaller~$k$,
\textsf{Clover} (without IS) is faster because the IS heuristic's
overhead exceeds its benefit on sparse subgraphs. On \textsf{com-LiveJournal}, \textsf{Clover+IS} finishes in under $10$ seconds
across $k = 3..12$, whereas \textsf{PivotScale} takes over $35$ hours at $k = 12$.
Relative speedups on other graphs fall in the range reported in the paper.

\noindent\emph{Expected reproduction time.}
End-to-end reproduction of the main comparison (all graphs, all binaries,
$k = 7$) takes approximately \textbf{8 hours} on a 192-thread Granite
Rapids node. The full $k = 3..12$ sweep (including all baselines) takes
several days, dominated by \textsf{PivotScale} on
\textsf{webbase-2001} at large~$k$. A downscaled variant running \textsf{Clover+IS} on
\textsf{com-LiveJournal} and \textsf{com-Orkut} completes in under
\textbf{30 minutes} on a modern workstation and suffices to exercise the
build and confirm the trends.

\noindent\emph{Artifact setup.}

\begin{enumerate}
    \item \textbf{Hardware dependencies.} x86-64 Linux node; the
    benchmark is multithreaded and scales to the available core count.
    The full sweep was run on a 192-thread Intel Granite Rapids node
    with 768~GB RAM. Any server with $\geq$128~GB RAM will run the full
    sweep; smaller machines handle the downscaled subset.
    \item \textbf{Software dependencies.} \texttt{gcc}~11 or newer with
    OpenMP and C++20 support; \texttt{make}; \texttt{curl},
    \texttt{tar}, \texttt{gunzip}; Python~3.8+ with \texttt{matplotlib}.
    \item \textbf{Data sets.} Five public graphs obtained at runtime by
    the script: \textsf{com-lj} and \textsf{com-orkut} from SNAP,
    \textsf{indochina-2004}, \textsf{uk-2005}, and \textsf{webbase-2001}
    from the SuiteSparse mirror of the LAW collection. No proprietary or
    locally-generated inputs.
    \item \textbf{Installation.}
\begin{verbatim}
git clone <repo-url>
cd CloverIS
\end{verbatim}
\end{enumerate}

\noindent\emph{Artifact execution.}
A single script, \texttt{reproduce.sh}, drives all four phases:
(i)~\texttt{make~-j all} builds \texttt{bin/clover}, \texttt{bin/clover\_is},
\texttt{bin/pivotscale}, and \texttt{bin/converter}; (ii)~each graph is
fetched, decompressed, and converted to the serialized \texttt{.sg} format
if not already present; (iii)~\textsf{PivotScale} is run once per $k$
while \textsf{Clover} and \textsf{Clover+IS} use a single invocation per
graph with a $k$-range option (\texttt{-c 3 -l 12}); (iv)~a small awk
parser pairs each \texttt{Counting Time:} line with its following
\texttt{k:} line and appends a row to \texttt{results.csv}.

\noindent\emph{Artifact analysis (including outputs).}
\texttt{plot.py} consumes \texttt{results.csv} and produces
\texttt{plots/runtime.\{png,svg\}}, a five-panel figure mirroring
Figure~5 in the paper. Reproduction of $C_1$ amounts to confirming that
the faster of \textsf{Clover} and \textsf{Clover+IS} outperforms
\textsf{PivotScale} at every~$k$ on every graph, and that the
speedup magnitudes fall within a small multiplicative tolerance of the
paper figures. Absolute timings depend on the reviewer's hardware.
Reproduction of $C_2$ corresponds to the \textsf{com-LiveJournal} panel:
\textsf{Clover+IS} completes the full $k$-sweep while \textsf{PivotScale}
exhibits super-linear growth and does not complete at large~$k$ within
the script's budget.

\subsection*{Artifact Evaluation (AE)}

\subsubsection*{A. Computational Artifact $A_1$}

\noindent\emph{Artifact setup.}
See the corresponding section in the Artifact Description. No additional
setup is required beyond cloning the repository.

\noindent\emph{Artifact execution.}

\noindent\textbf{Building.}
\begin{verbatim}
make -j all
\end{verbatim}
Produces four binaries in \texttt{bin/}: \texttt{pivotscale},
\texttt{clover}, \texttt{clover\_is}, and \texttt{converter}. The first
three compile from two independent source trees
(\texttt{src/pivotscale/} and \texttt{src/clover/}); \textsf{Clover} and
\textsf{Clover+IS} share a single source file, with the IS path gated
behind a compile-time flag (\texttt{-DENABLE\_IS=0} or~\texttt{=1}) so
that the two binaries contain disjoint executable code.

\noindent\textbf{Running.}
\begin{verbatim}
./reproduce.sh
\end{verbatim}
All remaining phases --- graph fetch, conversion, benchmark sweep,
\texttt{results.csv} assembly, and plotting --- run to completion
without further intervention. Intermediate downloads are deleted after
conversion; the \texttt{.sg} files are retained so that a re-run skips
the fetch phase. The script is idempotent with respect to existing
\texttt{.sg} inputs.

\noindent\textbf{Collecting and plotting results.}
\texttt{reproduce.sh} writes \texttt{results.csv} (one row per
\((graph, binary, k)\)) and then invokes \texttt{python3 plot.py
results.csv plots/}, which writes \texttt{plots/runtime.png} and
\texttt{plots/runtime.svg}.

\noindent\emph{Artifact analysis (including outputs).}
The reproduction is considered successful when the produced
\texttt{runtime} figure exhibits the same qualitative shape as the paper
figure --- \textsf{Clover+IS} fastest, \textsf{PivotScale} slowest, both
separated by one to four orders of magnitude at large~$k$ --- and when
the relative speedups in \texttt{results.csv} fall within a small
tolerance of the corresponding entries in the shipped
\texttt{paper\_results.csv}. Exact bitwise or timing equality is not
expected and, per the SC26 guidelines, is not the goal of the
\textsc{Results Reproduced} badge.


\section*{Appendix: Artifact Description/Artifact Evaluation}

\subsection*{Artifact Description (AD)}

\subsubsection*{I. Overview of Contributions and Artifacts}

\paragraph{A. Paper's Main Contributions}
The paper introduces \textsf{Clover}, a parallel $k$-clique counting
algorithm based on a Cover-oriented pivoting search, and its extension
\textsf{Clover+IS}, which counts independent sets on the complement graph
when a residual subproblem becomes dense. In the context of the artifact
description, the relevant contributions are:

\begin{description}
    \item[$C_1$] \textsf{Clover} and \textsf{Clover+IS} outperform the
    state-of-the-art \textsf{PivotScale} baseline by one to four orders of
    magnitude for large cliques on large real-world graphs.
    \item[$C_2$] \textsf{Clover+IS} is the first system to count
    $k$-cliques of arbitrary $k$ on \textsf{com-LiveJournal}, a problem
    that was previously intractable.
\end{description}

\paragraph{B. Computational Artifacts}
\begin{center}
\begin{tabular}{lll}
\toprule
Artifact ID & Contributions & Related Paper Elements \\
\midrule
$A_1$ & $C_1$, $C_2$ & Figure~5, Table~II \\
\bottomrule
\end{tabular}
\end{center}

$A_1$: \url{https://github.com/<TBD>/CloverIS} (final DOI assigned via
Zenodo by the Artifact Freeze deadline).

\subsubsection*{II. Artifact Identification}

\paragraph{A. Computational Artifact $A_1$}

\noindent\emph{Relation to contributions.}
$A_1$ contains the full source of \textsf{Clover}, \textsf{Clover+IS}, and
the \textsf{PivotScale} baseline, together with the download/convert
pipeline for the five input graphs and the plotting driver. Running the
single entry-point script produces the per-$(graph, k)$ counting times
that substantiate $C_1$ and $C_2$.

\noindent\emph{Expected results.}
The faster of \textsf{Clover} and \textsf{Clover+IS} outperforms
\textsf{PivotScale} on every graph at every~$k$ tested.
\textsf{Clover+IS} is fastest for $k \geq 5$; at smaller~$k$,
\textsf{Clover} (without IS) is faster because the IS heuristic's
overhead exceeds its benefit on sparse subgraphs. On \textsf{com-LiveJournal}, \textsf{Clover+IS} finishes in under $10$ seconds
across $k = 3..12$, whereas \textsf{PivotScale} takes over $35$ hours at $k = 12$.
Relative speedups on other graphs fall in the range reported in the paper.

\noindent\emph{Expected reproduction time.}
End-to-end reproduction of the main comparison (all graphs, all binaries,
$k = 7$) takes approximately \textbf{8 hours} on a 192-thread Granite
Rapids node. The full $k = 3..12$ sweep (including all baselines) takes
several days, dominated by \textsf{PivotScale} on
\textsf{webbase-2001} at large~$k$. A downscaled variant running \textsf{Clover+IS} on
\textsf{com-LiveJournal} and \textsf{com-Orkut} completes in under
\textbf{30 minutes} on a modern workstation and suffices to exercise the
build and confirm the trends.

\noindent\emph{Artifact setup.}

\begin{enumerate}
    \item \textbf{Hardware dependencies.} x86-64 Linux node; the
    benchmark is multithreaded and scales to the available core count.
    The full sweep was run on a 192-thread Intel Granite Rapids node
    with 768~GB RAM. Any server with $\geq$128~GB RAM will run the full
    sweep; smaller machines handle the downscaled subset.
    \item \textbf{Software dependencies.} \texttt{gcc}~11 or newer with
    OpenMP and C++20 support; \texttt{make}; \texttt{curl},
    \texttt{tar}, \texttt{gunzip}; Python~3.8+ with \texttt{matplotlib}.
    \item \textbf{Data sets.} Five public graphs obtained at runtime by
    the script: \textsf{com-lj} and \textsf{com-orkut} from SNAP,
    \textsf{indochina-2004}, \textsf{uk-2005}, and \textsf{webbase-2001}
    from the SuiteSparse mirror of the LAW collection. No proprietary or
    locally-generated inputs.
    \item \textbf{Installation.}
\begin{verbatim}
git clone <repo-url>
cd CloverIS
\end{verbatim}
\end{enumerate}

\noindent\emph{Artifact execution.}
A single script, \texttt{reproduce.sh}, drives all four phases:
(i)~\texttt{make~-j all} builds \texttt{bin/clover}, \texttt{bin/clover\_is},
\texttt{bin/pivotscale}, and \texttt{bin/converter}; (ii)~each graph is
fetched, decompressed, and converted to the serialized \texttt{.sg} format
if not already present; (iii)~\textsf{PivotScale} is run once per $k$
while \textsf{Clover} and \textsf{Clover+IS} use a single invocation per
graph with a $k$-range option (\texttt{-c 3 -l 12}); (iv)~a small awk
parser pairs each \texttt{Counting Time:} line with its following
\texttt{k:} line and appends a row to \texttt{results.csv}.

\noindent\emph{Artifact analysis (including outputs).}
\texttt{plot.py} consumes \texttt{results.csv} and produces
\texttt{plots/runtime.\{png,svg\}}, a five-panel figure mirroring
Figure~5 in the paper. Reproduction of $C_1$ amounts to confirming that
the faster of \textsf{Clover} and \textsf{Clover+IS} outperforms
\textsf{PivotScale} at every~$k$ on every graph, and that the
speedup magnitudes fall within a small multiplicative tolerance of the
paper figures. Absolute timings depend on the reviewer's hardware.
Reproduction of $C_2$ corresponds to the \textsf{com-LiveJournal} panel:
\textsf{Clover+IS} completes the full $k$-sweep while \textsf{PivotScale}
exhibits super-linear growth and does not complete at large~$k$ within
the script's budget.

\subsection*{Artifact Evaluation (AE)}

\subsubsection*{A. Computational Artifact $A_1$}

\noindent\emph{Artifact setup.}
See the corresponding section in the Artifact Description. No additional
setup is required beyond cloning the repository.

\noindent\emph{Artifact execution.}

\noindent\textbf{Building.}
\begin{verbatim}
make -j all
\end{verbatim}
Produces four binaries in \texttt{bin/}: \texttt{pivotscale},
\texttt{clover}, \texttt{clover\_is}, and \texttt{converter}. The first
three compile from two independent source trees
(\texttt{src/pivotscale/} and \texttt{src/clover/}); \textsf{Clover} and
\textsf{Clover+IS} share a single source file, with the IS path gated
behind a compile-time flag (\texttt{-DENABLE\_IS=0} or~\texttt{=1}) so
that the two binaries contain disjoint executable code.

\noindent\textbf{Running.}
\begin{verbatim}
./reproduce.sh
\end{verbatim}
All remaining phases --- graph fetch, conversion, benchmark sweep,
\texttt{results.csv} assembly, and plotting --- run to completion
without further intervention. Intermediate downloads are deleted after
conversion; the \texttt{.sg} files are retained so that a re-run skips
the fetch phase. The script is idempotent with respect to existing
\texttt{.sg} inputs.

\noindent\textbf{Collecting and plotting results.}
\texttt{reproduce.sh} writes \texttt{results.csv} (one row per
\((graph, binary, k)\)) and then invokes \texttt{python3 plot.py
results.csv plots/}, which writes \texttt{plots/runtime.png} and
\texttt{plots/runtime.svg}.

\noindent\emph{Artifact analysis (including outputs).}
The reproduction is considered successful when the produced
\texttt{runtime} figure exhibits the same qualitative shape as the paper
figure --- \textsf{Clover+IS} fastest, \textsf{PivotScale} slowest, both
separated by one to four orders of magnitude at large~$k$ --- and when
the relative speedups in \texttt{results.csv} fall within a small
tolerance of the corresponding entries in the shipped
\texttt{paper\_results.csv}. Exact bitwise or timing equality is not
expected and, per the SC26 guidelines, is not the goal of the
\textsc{Results Reproduced} badge.


\section*{Appendix: Artifact Description/Artifact Evaluation}

\subsection*{Artifact Description (AD)}

\subsubsection*{I. Overview of Contributions and Artifacts}

\paragraph{A. Paper's Main Contributions}
The paper introduces \textsf{Clover}, a parallel $k$-clique counting
algorithm based on a Cover-oriented pivoting search, and its extension
\textsf{Clover+IS}, which counts independent sets on the complement graph
when a residual subproblem becomes dense. In the context of the artifact
description, the relevant contributions are:

\begin{description}
    \item[$C_1$] \textsf{Clover} and \textsf{Clover+IS} outperform the
    state-of-the-art \textsf{PivotScale} baseline by one to four orders of
    magnitude for large cliques on large real-world graphs.
    \item[$C_2$] \textsf{Clover+IS} is the first system to count
    $k$-cliques of arbitrary $k$ on \textsf{com-LiveJournal}, a problem
    that was previously intractable.
\end{description}

\paragraph{B. Computational Artifacts}
\begin{center}
\begin{tabular}{lll}
\toprule
Artifact ID & Contributions & Related Paper Elements \\
\midrule
$A_1$ & $C_1$, $C_2$ & Figure~5, Table~II \\
\bottomrule
\end{tabular}
\end{center}

$A_1$: \url{https://github.com/<TBD>/CloverIS} (final DOI assigned via
Zenodo by the Artifact Freeze deadline).

\subsubsection*{II. Artifact Identification}

\paragraph{A. Computational Artifact $A_1$}

\noindent\emph{Relation to contributions.}
$A_1$ contains the full source of \textsf{Clover}, \textsf{Clover+IS}, and
the \textsf{PivotScale} baseline, together with the download/convert
pipeline for the five input graphs and the plotting driver. Running the
single entry-point script produces the per-$(graph, k)$ counting times
that substantiate $C_1$ and $C_2$.

\noindent\emph{Expected results.}
The faster of \textsf{Clover} and \textsf{Clover+IS} outperforms
\textsf{PivotScale} on every graph at every~$k$ tested.
\textsf{Clover+IS} is fastest for $k \geq 5$; at smaller~$k$,
\textsf{Clover} (without IS) is faster because the IS heuristic's
overhead exceeds its benefit on sparse subgraphs. On \textsf{com-LiveJournal}, \textsf{Clover+IS} finishes in under $10$ seconds
across $k = 3..12$, whereas \textsf{PivotScale} takes over $35$ hours at $k = 12$.
Relative speedups on other graphs fall in the range reported in the paper.

\noindent\emph{Expected reproduction time.}
End-to-end reproduction of the main comparison (all graphs, all binaries,
$k = 7$) takes approximately \textbf{8 hours} on a 192-thread Granite
Rapids node. The full $k = 3..12$ sweep (including all baselines) takes
several days, dominated by \textsf{PivotScale} on
\textsf{webbase-2001} at large~$k$. A downscaled variant running \textsf{Clover+IS} on
\textsf{com-LiveJournal} and \textsf{com-Orkut} completes in under
\textbf{30 minutes} on a modern workstation and suffices to exercise the
build and confirm the trends.

\noindent\emph{Artifact setup.}

\begin{enumerate}
    \item \textbf{Hardware dependencies.} x86-64 Linux node; the
    benchmark is multithreaded and scales to the available core count.
    The full sweep was run on a 192-thread Intel Granite Rapids node
    with 768~GB RAM. Any server with $\geq$128~GB RAM will run the full
    sweep; smaller machines handle the downscaled subset.
    \item \textbf{Software dependencies.} \texttt{gcc}~11 or newer with
    OpenMP and C++20 support; \texttt{make}; \texttt{curl},
    \texttt{tar}, \texttt{gunzip}; Python~3.8+ with \texttt{matplotlib}.
    \item \textbf{Data sets.} Five public graphs obtained at runtime by
    the script: \textsf{com-lj} and \textsf{com-orkut} from SNAP,
    \textsf{indochina-2004}, \textsf{uk-2005}, and \textsf{webbase-2001}
    from the SuiteSparse mirror of the LAW collection. No proprietary or
    locally-generated inputs.
    \item \textbf{Installation.}
\begin{verbatim}
git clone <repo-url>
cd CloverIS
\end{verbatim}
\end{enumerate}

\noindent\emph{Artifact execution.}
A single script, \texttt{reproduce.sh}, drives all four phases:
(i)~\texttt{make~-j all} builds \texttt{bin/clover}, \texttt{bin/clover\_is},
\texttt{bin/pivotscale}, and \texttt{bin/converter}; (ii)~each graph is
fetched, decompressed, and converted to the serialized \texttt{.sg} format
if not already present; (iii)~\textsf{PivotScale} is run once per $k$
while \textsf{Clover} and \textsf{Clover+IS} use a single invocation per
graph with a $k$-range option (\texttt{-c 3 -l 12}); (iv)~a small awk
parser pairs each \texttt{Counting Time:} line with its following
\texttt{k:} line and appends a row to \texttt{results.csv}.

\noindent\emph{Artifact analysis (including outputs).}
\texttt{plot.py} consumes \texttt{results.csv} and produces
\texttt{plots/runtime.\{png,svg\}}, a five-panel figure mirroring
Figure~5 in the paper. Reproduction of $C_1$ amounts to confirming that
the faster of \textsf{Clover} and \textsf{Clover+IS} outperforms
\textsf{PivotScale} at every~$k$ on every graph, and that the
speedup magnitudes fall within a small multiplicative tolerance of the
paper figures. Absolute timings depend on the reviewer's hardware.
Reproduction of $C_2$ corresponds to the \textsf{com-LiveJournal} panel:
\textsf{Clover+IS} completes the full $k$-sweep while \textsf{PivotScale}
exhibits super-linear growth and does not complete at large~$k$ within
the script's budget.

\subsection*{Artifact Evaluation (AE)}

\subsubsection*{A. Computational Artifact $A_1$}

\noindent\emph{Artifact setup.}
See the corresponding section in the Artifact Description. No additional
setup is required beyond cloning the repository.

\noindent\emph{Artifact execution.}

\noindent\textbf{Building.}
\begin{verbatim}
make -j all
\end{verbatim}
Produces four binaries in \texttt{bin/}: \texttt{pivotscale},
\texttt{clover}, \texttt{clover\_is}, and \texttt{converter}. The first
three compile from two independent source trees
(\texttt{src/pivotscale/} and \texttt{src/clover/}); \textsf{Clover} and
\textsf{Clover+IS} share a single source file, with the IS path gated
behind a compile-time flag (\texttt{-DENABLE\_IS=0} or~\texttt{=1}) so
that the two binaries contain disjoint executable code.

\noindent\textbf{Running.}
\begin{verbatim}
./reproduce.sh
\end{verbatim}
All remaining phases --- graph fetch, conversion, benchmark sweep,
\texttt{results.csv} assembly, and plotting --- run to completion
without further intervention. Intermediate downloads are deleted after
conversion; the \texttt{.sg} files are retained so that a re-run skips
the fetch phase. The script is idempotent with respect to existing
\texttt{.sg} inputs.

\noindent\textbf{Collecting and plotting results.}
\texttt{reproduce.sh} writes \texttt{results.csv} (one row per
\((graph, binary, k)\)) and then invokes \texttt{python3 plot.py
results.csv plots/}, which writes \texttt{plots/runtime.png} and
\texttt{plots/runtime.svg}.

\noindent\emph{Artifact analysis (including outputs).}
The reproduction is considered successful when the produced
\texttt{runtime} figure exhibits the same qualitative shape as the paper
figure --- \textsf{Clover+IS} fastest, \textsf{PivotScale} slowest, both
separated by one to four orders of magnitude at large~$k$ --- and when
the relative speedups in \texttt{results.csv} fall within a small
tolerance of the corresponding entries in the shipped
\texttt{paper\_results.csv}. Exact bitwise or timing equality is not
expected and, per the SC26 guidelines, is not the goal of the
\textsc{Results Reproduced} badge.


\section*{Appendix: Artifact Description/Artifact Evaluation}

\subsection*{Artifact Description (AD)}

\subsubsection*{I. Overview of Contributions and Artifacts}

\paragraph{A. Paper's Main Contributions}
The paper introduces \textsf{Clover}, a parallel $k$-clique counting
algorithm based on a Cover-oriented pivoting search, and its extension
\textsf{Clover+IS}, which counts independent sets on the complement graph
when a residual subproblem becomes dense. In the context of the artifact
description, the relevant contributions are:

\begin{description}
    \item[$C_1$] \textsf{Clover} and \textsf{Clover+IS} outperform the
    state-of-the-art \textsf{PivotScale} baseline by one to four orders of
    magnitude for large cliques on large real-world graphs.
    \item[$C_2$] \textsf{Clover+IS} is the first system to count
    $k$-cliques of arbitrary $k$ on \textsf{com-LiveJournal}, a problem
    that was previously intractable.
\end{description}

\paragraph{B. Computational Artifacts}
\begin{center}
\begin{tabular}{lll}
\toprule
Artifact ID & Contributions & Related Paper Elements \\
\midrule
$A_1$ & $C_1$, $C_2$ & Figure~5, Table~II \\
\bottomrule
\end{tabular}
\end{center}

$A_1$: \url{https://github.com/<TBD>/CloverIS} (final DOI assigned via
Zenodo by the Artifact Freeze deadline).

\subsubsection*{II. Artifact Identification}

\paragraph{A. Computational Artifact $A_1$}

\noindent\emph{Relation to contributions.}
$A_1$ contains the full source of \textsf{Clover}, \textsf{Clover+IS}, and
the \textsf{PivotScale} baseline, together with the download/convert
pipeline for the five input graphs and the plotting driver. Running the
single entry-point script produces the per-$(graph, k)$ counting times
that substantiate $C_1$ and $C_2$.

\noindent\emph{Expected results.}
The faster of \textsf{Clover} and \textsf{Clover+IS} outperforms
\textsf{PivotScale} on every graph at every~$k$ tested.
\textsf{Clover+IS} is fastest for $k \geq 5$; at smaller~$k$,
\textsf{Clover} (without IS) is faster because the IS heuristic's
overhead exceeds its benefit on sparse subgraphs. On \textsf{com-LiveJournal}, \textsf{Clover+IS} finishes in under $10$ seconds
across $k = 3..12$, whereas \textsf{PivotScale} takes over $35$ hours at $k = 12$.
Relative speedups on other graphs fall in the range reported in the paper.

\noindent\emph{Expected reproduction time.}
End-to-end reproduction of the main comparison (all graphs, all binaries,
$k = 7$) takes approximately \textbf{8 hours} on a 192-thread Granite
Rapids node. The full $k = 3..12$ sweep (including all baselines) takes
several days, dominated by \textsf{PivotScale} on
\textsf{webbase-2001} at large~$k$. A downscaled variant running \textsf{Clover+IS} on
\textsf{com-LiveJournal} and \textsf{com-Orkut} completes in under
\textbf{30 minutes} on a modern workstation and suffices to exercise the
build and confirm the trends.

\noindent\emph{Artifact setup.}

\begin{enumerate}
    \item \textbf{Hardware dependencies.} x86-64 Linux node; the
    benchmark is multithreaded and scales to the available core count.
    The full sweep was run on a 192-thread Intel Granite Rapids node
    with 768~GB RAM. Any server with $\geq$128~GB RAM will run the full
    sweep; smaller machines handle the downscaled subset.
    \item \textbf{Software dependencies.} \texttt{gcc}~11 or newer with
    OpenMP and C++20 support; \texttt{make}; \texttt{curl},
    \texttt{tar}, \texttt{gunzip}; Python~3.8+ with \texttt{matplotlib}.
    \item \textbf{Data sets.} Five public graphs obtained at runtime by
    the script: \textsf{com-lj} and \textsf{com-orkut} from SNAP,
    \textsf{indochina-2004}, \textsf{uk-2005}, and \textsf{webbase-2001}
    from the SuiteSparse mirror of the LAW collection. No proprietary or
    locally-generated inputs.
    \item \textbf{Installation.}
\begin{verbatim}
git clone <repo-url>
cd CloverIS
\end{verbatim}
\end{enumerate}

\noindent\emph{Artifact execution.}
A single script, \texttt{reproduce.sh}, drives all four phases:
(i)~\texttt{make~-j all} builds \texttt{bin/clover}, \texttt{bin/clover\_is},
\texttt{bin/pivotscale}, and \texttt{bin/converter}; (ii)~each graph is
fetched, decompressed, and converted to the serialized \texttt{.sg} format
if not already present; (iii)~\textsf{PivotScale} is run once per $k$
while \textsf{Clover} and \textsf{Clover+IS} use a single invocation per
graph with a $k$-range option (\texttt{-c 3 -l 12}); (iv)~a small awk
parser pairs each \texttt{Counting Time:} line with its following
\texttt{k:} line and appends a row to \texttt{results.csv}.

\noindent\emph{Artifact analysis (including outputs).}
\texttt{plot.py} consumes \texttt{results.csv} and produces
\texttt{plots/runtime.\{png,svg\}}, a five-panel figure mirroring
Figure~5 in the paper. Reproduction of $C_1$ amounts to confirming that
the faster of \textsf{Clover} and \textsf{Clover+IS} outperforms
\textsf{PivotScale} at every~$k$ on every graph, and that the
speedup magnitudes fall within a small multiplicative tolerance of the
paper figures. Absolute timings depend on the reviewer's hardware.
Reproduction of $C_2$ corresponds to the \textsf{com-LiveJournal} panel:
\textsf{Clover+IS} completes the full $k$-sweep while \textsf{PivotScale}
exhibits super-linear growth and does not complete at large~$k$ within
the script's budget.

\subsection*{Artifact Evaluation (AE)}

\subsubsection*{A. Computational Artifact $A_1$}

\noindent\emph{Artifact setup.}
See the corresponding section in the Artifact Description. No additional
setup is required beyond cloning the repository.

\noindent\emph{Artifact execution.}

\noindent\textbf{Building.}
\begin{verbatim}
make -j all
\end{verbatim}
Produces four binaries in \texttt{bin/}: \texttt{pivotscale},
\texttt{clover}, \texttt{clover\_is}, and \texttt{converter}. The first
three compile from two independent source trees
(\texttt{src/pivotscale/} and \texttt{src/clover/}); \textsf{Clover} and
\textsf{Clover+IS} share a single source file, with the IS path gated
behind a compile-time flag (\texttt{-DENABLE\_IS=0} or~\texttt{=1}) so
that the two binaries contain disjoint executable code.

\noindent\textbf{Running.}
\begin{verbatim}
./reproduce.sh
\end{verbatim}
All remaining phases --- graph fetch, conversion, benchmark sweep,
\texttt{results.csv} assembly, and plotting --- run to completion
without further intervention. Intermediate downloads are deleted after
conversion; the \texttt{.sg} files are retained so that a re-run skips
the fetch phase. The script is idempotent with respect to existing
\texttt{.sg} inputs.

\noindent\textbf{Collecting and plotting results.}
\texttt{reproduce.sh} writes \texttt{results.csv} (one row per
\((graph, binary, k)\)) and then invokes \texttt{python3 plot.py
results.csv plots/}, which writes \texttt{plots/runtime.png} and
\texttt{plots/runtime.svg}.

\noindent\emph{Artifact analysis (including outputs).}
The reproduction is considered successful when the produced
\texttt{runtime} figure exhibits the same qualitative shape as the paper
figure --- \textsf{Clover+IS} fastest, \textsf{PivotScale} slowest, both
separated by one to four orders of magnitude at large~$k$ --- and when
the relative speedups in \texttt{results.csv} fall within a small
tolerance of the corresponding entries in the shipped
\texttt{paper\_results.csv}. Exact bitwise or timing equality is not
expected and, per the SC26 guidelines, is not the goal of the
\textsc{Results Reproduced} badge.
